\chapter{Down Syndrome}

\section*{Definition and Overview}
Down syndrome is a genetic condition caused by the presence of an extra copy of chromosome 21. It affects development, causing characteristic physical features, intellectual disability of varying degree, and an increased risk of certain health conditions, including heart defects, thyroid problems, and hearing issues. Down syndrome occurs in people of all backgrounds and is not caused by anything parents did or did not do. With good health care, early intervention, education, and family support, most people with Down syndrome live healthy, active, and fulfilling lives. This chapter explains the causes, features, and care of Down syndrome.

\section*{Epidemiology}
Down syndrome is the most common chromosomal condition, occurring in around 1 in 800--1,000 live births. The risk increases with maternal age, from about 1 in 1,500 in a 20-year-old mother to around 1 in 100 at age 40, but because most babies are born to younger mothers, most children with Down syndrome are born to women under 35. In Australia, around 250--300 babies with Down syndrome are born each year. Advances in medical and educational support have greatly improved life expectancy, with most people living into their 60s or beyond.

\section*{Aetiology and Pathophysiology}
\begin{itemize}
    \item \textbf{Trisomy 21:} most people with Down syndrome have three copies of chromosome 21 instead of two
    \item \textbf{Translocation:} in a small proportion, part of chromosome 21 is attached to another chromosome
    \item \textbf{Mosaicism:} in a minority, only some cells have the extra chromosome, with milder features
    \item \textbf{Extra genetic material:} the extra chromosome affects development of the brain, face, and organs
    \item \textbf{Intellectual disability:} usually mild to moderate, with strengths in social skills
    \item \textbf{Associated conditions:} heart defects, gastrointestinal problems, hypothyroidism, and immune issues are more common
\end{itemize}

\section*{Clinical Features}
\begin{itemize}
    \item Characteristic facial features: flat facial profile, upward-slanting eyes, and a small mouth
    \item Low muscle tone (hypotonia) in infancy
    \item Small stature and a single palmar crease in many
    \item Delayed motor and speech development
    \item Mild to moderate intellectual disability with individual variation
    \item Increased risk of congenital heart defects, hearing and vision problems, and thyroid disease
    \item Warm, sociable personality is common but not universal
\end{itemize}

\section*{Differential Diagnosis}
\begin{itemize}
    \item Other chromosomal and genetic conditions with overlapping features
    \item Other causes of intellectual disability and developmental delay
    \item Normal variation in appearance and development
    \item Fragile X syndrome and other genetic syndromes
\end{itemize}

\section*{When to Seek Medical Care}
Prenatal screening and diagnostic testing are offered to all pregnant women. After birth, children with Down syndrome need regular medical reviews to detect and manage associated conditions early. Parents should seek advice about early intervention services and supports.

\textbf{Call 000 (or local emergency services) for:}
\begin{itemize}
    \item Difficulty breathing, blue lips, or signs of heart problems in a baby
    \item A baby who is very floppy, unresponsive, or not feeding
    \item Sudden neurological symptoms or seizures
    \item Any serious illness or injury
\end{itemize}

\section*{Diagnosis and Investigations}
\begin{itemize}
    \item \textbf{Prenatal screening:} combined first-trimester screening and non-invasive prenatal testing (NIPT) estimate the risk
    \item \textbf{Prenatal diagnosis:} chorionic villus sampling or amniocentesis confirm the diagnosis
    \item \textbf{Postnatal diagnosis:} blood tests (karyotype) confirm the condition after birth
    \item \textbf{Heart assessment:} echocardiogram in the newborn period
    \item \textbf{Hearing and vision tests:} regular assessment in childhood
    \item \textbf{Thyroid and other tests:} screening for associated conditions
\end{itemize}

\section*{Management}
\begin{itemize}
    \item \textbf{Early intervention:} physiotherapy, speech therapy, and occupational therapy from infancy
    \item \textbf{Education support:} inclusive schooling with individualised support
    \item \textbf{Treatment of heart defects:} surgery for congenital heart disease when needed
    \item \textbf{Thyroid treatment:} medication for hypothyroidism
    \item \textbf{Hearing and vision care:} treatment of hearing loss and vision problems
    \item \textbf{Health surveillance:} regular reviews of growth, development, and associated conditions
    \item \textbf{Family support:} information, counselling, and connection with support organisations
    \item \textbf{Adult services:} supported employment, independent living, and social programs
\end{itemize}

\section*{Complications}
\begin{itemize}
    \item Congenital heart disease, the most significant health concern
    \item Hearing loss, vision problems, and recurrent infections
    \item Hypothyroidism and other endocrine conditions
    \item Gastrointestinal conditions, including coeliac disease
    \item Obesity and sleep apnoea
    \item Early-onset Alzheimer's disease in older adults
    \item Leukaemia, which is more common but still rare
\end{itemize}

\section*{Prognosis and Outcomes}
The outlook for people with Down syndrome has improved dramatically. Most children with Down syndrome are healthy, attend school, and participate in community life, and life expectancy is now around 60 years. Early intervention and good medical care maximise development and independence. Many adults with Down syndrome live semi-independently, work, and enjoy rich social lives. With love, support, and appropriate services, people with Down syndrome thrive.

\section*{Prevention and Self-Care}
\begin{itemize}
    \item Attend prenatal screening if you wish to know your risk
    \item Connect early with Down syndrome support organisations
    \item Access early intervention services promptly
    \item Attend all recommended medical and developmental reviews
    \item Ensure hearing, vision, and thyroid checks are done
    \item Support healthy eating, activity, and weight
    \item Plan for education, employment, and adult life
    \item Look after the wellbeing of siblings and parents
\end{itemize}

\section*{Sources and Further Reading}
The following reputable sources provide further information for healthcare students and patients seeking reliable, current advice about Down syndrome.

\begin{itemize}
    \item Down Syndrome Australia. \textit{Information and support for families.} https://www.downsyndrome.org.au/
    \item Down Syndrome NSW. \textit{Resources and support.} https://www.downsydney.org/
    \item Healthdirect Australia. \textit{Down syndrome: information.} https://www.healthdirect.gov.au/
    \item National Down Syndrome Society. \textit{Information about Down syndrome.} https://ndss.org/
    \item MedlinePlus. \textit{Down syndrome.} https://medlineplus.gov/
    \item Raising Children Network. \textit{Down syndrome: parenting and development.} https://raisingchildren.net.au/
\end{itemize}
